\documentclass{amsart}

% 数学相关宏包
\usepackage{amsmath,mathtools,amsthm,amssymb}
\usepackage{bbm}
\usepackage{mathrsfs}
\usepackage{centernot}

% 图形和绘图
\usepackage{graphicx,float}
\usepackage{tikz,tikz-cd}
\usepackage{pgfplots}
\usepackage[framemethod=TikZ]{mdframed}
\usetikzlibrary{decorations.pathmorphing,positioning,arrows,automata,calc,shapes,patterns}
\pgfplotsset{compat=1.18}
\usepgfplotslibrary{statistics}

% 表格和列表
\usepackage{booktabs,multirow,colortbl}
\usepackage{enumitem}
\setlist[itemize,enumerate,description]{noitemsep}
\setlistdepth{9}

% 算法
\usepackage{algorithm,algorithmic}

% 字体和颜色
\usepackage{xcolor,listings}
\usepackage{xeCJK}
\setCJKmainfont{Noto Serif CJK SC}

% 页面和链接
\usepackage{fancyhdr,caption}
\usepackage{hyperref}
\usepackage{cleveref}

% 工具宏包
\usepackage{etoolbox,letltxmacro,docmute}
\usepackage{gensymb}

% 自定义设置
\renewcommand{\arraystretch}{1.5}
\let\originalmathbb\mathbb
\renewcommand{\mathbb}[1]{%
    \ifstrequal{#1}{1}{\mathbbm{#1}}{\originalmathbb{#1}}%
}

% 颜色定义
\definecolor{lightblue}{RGB}{173,216,230}
\definecolor{lightgreen}{RGB}{144,238,144}
\definecolor{lightyellow}{RGB}{255,255,224}
\definecolor{lightgray}{RGB}{211,211,211}

		% 超链接设置
		\hypersetup{
		    colorlinks=true,
		    linkcolor=red,
		    filecolor=magenta,
		    urlcolor=cyan,
		    pdftitle={\@title},
		    pdfauthor={\@author},
		    pdfsubject={数学笔记},
		    pdfkeywords={数学},
		    bookmarksnumbered=true,
		    bookmarksopen=true,
		    bookmarksopenlevel=6,
		    pdfstartview=FitH,
		    pdfpagemode=UseOutlines,
		    pdfborder={0 0 0}
		}

% 定理环境
\theoremstyle{plain}
\newtheorem{theorem}{Theorem}[section]
\newtheorem{lemma}[theorem]{Lemma}
\newtheorem{corollary}[theorem]{Corollary}
\newtheorem{proposition}[theorem]{Proposition}

\theoremstyle{definition}
\newtheorem{definition}[theorem]{Definition}
\newtheorem{example}[theorem]{Example}
\newtheorem{exercise}[theorem]{Exercise}
\newtheorem{problem}[theorem]{Problem}

\theoremstyle{remark}
\newtheorem{remark}[theorem]{Remark}
\newtheorem{note}[theorem]{Note}
\newtheorem{observation}[theorem]{Observation}

% 解答环境
\newtheorem*{solution}{Solution}
\newtheorem*{proof*}{Proof}

% 图片路径
\graphicspath{{F:/资料/2025-summer/figures/}}

\author{乐绎华, Yue Yihua}
\address{中山大学, Sun Yat-sen University}
\email{yueyh@mail2.sysu.edu.cn}
\date{\today}
\title{Mathematical Notes}

\begin{document}

\maketitle
\tableofcontents

\section{曲面论基本定理}

\subsection{Christoffel 记号，Gauss 公式，Wiengarten 公式}

\begin{figure}[H]
\centering
\includegraphics[width=\textwidth]{differential-geometry-2025070923.png}
% \caption{}
\label{}
\end{figure}

对于 $r(u^{1},v^{1})$，
\[
\mathrm{d}r=\sum_{\alpha=1}^{2} r_{\alpha}\mathrm{d}u^{\alpha}=r_{\alpha}\mathrm{d}u^{\alpha}
\]
第一基本形式：
\[
\mathrm{I}=\lvert \mathrm{d}r \rvert ^{2}=\sum_{\alpha,\beta=1}^{2} \underbrace{ (r_{\alpha}\cdot r_{\beta}) }_{ \eqqcolon g_{\alpha\beta} }\mathrm{d}u^{\alpha}\mathrm{d}u^{\beta}=g_{\alpha\beta}\mathrm{d}u^{\alpha}\mathrm{d}u^{\beta}
\]
第二基本形式：
\[
\mathrm{II}=\sum_{\alpha,\beta=1}^{2} \underbrace{ (r_{\alpha\beta}\cdot n) }_{ \eqqcolon  b_{\alpha\beta} }\mathrm{d}u^{\alpha}\mathrm{d}u^{\beta}=b_{\alpha\beta}\mathrm{d}u^{\alpha}\mathrm{d}u^{\beta}
\]
$b_{\alpha\beta}=r_{\alpha\beta}\cdot n=-r_{\alpha}\cdot n_{\beta}=-r_{\beta}\cdot n_{\alpha}$. 记 $g\coloneqq \det(g_{\alpha\beta}),b\coloneqq \det(b_{\alpha\beta})$, $(g_{\alpha\beta})$ 显然正定，其逆记作 $(g^{\alpha\beta})$. 有如下关系
\[
g^{\alpha\beta}g_{\beta\gamma}=\delta_{\gamma}^{\alpha}
\]
\[
\begin{pmatrix}
g^{11}  & g^{12 } \\
g^{21} & g^{22}
\end{pmatrix}=\frac{1}{g}\begin{pmatrix}
g_{22} & -g_{12} \\
-g_{21} & g_{11}
\end{pmatrix}
\]
曲面 $S:r=r(u^{1},u^{2})$ 在每点处有自然标架 $\{ r;r_1,r_2,n \}$，其中 $r_1,r_2$ 由 $\partial_1r,\partial_2r$ 定义，$n\coloneqq\frac{r_1\times r_2}{\lvert r_1\times r_2 \rvert}$. 点在曲面上运动的时候，对应的自然标架也是运动的，考虑描述其变化的方程：
\begin{equation}
\begin{cases}
\frac{ \partial r_{\alpha} }{ \partial u^{\beta} }  & =\Gamma^{\gamma}_{\alpha\beta}r_{\gamma}+C_{\alpha\beta}n \\
\frac{ \partial n }{ \partial u^{\beta} }  & =D^{\gamma}_{\beta}r_{\gamma }+D_{\beta}n
\end{cases}
\label{1233fd}
\end{equation}

其中 $\Gamma^{\gamma}_{\alpha\beta},C_{\alpha\beta},D^{\gamma}_{\beta},D_{\beta}$ 都是待定系数. 分别用 $r_{\xi},n$ 与上式做内积得到
\[
D_{\beta}=n_{\beta}\cdot n=-n\cdot n_{\beta}=0
\]
\[
C_{\alpha\beta}=r_{\alpha\beta}\cdot n=b_{\alpha\beta}
\]
\begin{equation}
D^{\gamma}_{\beta}g_{\gamma \xi}=D^{\gamma}_{\beta}r_{\gamma}\cdot r_{\gamma}=n_{\beta}\cdot r_{\xi}=-b_{\beta \xi}
\label{cb4036}
\end{equation}

于是
\begin{equation}
D^{\gamma}_{\beta}=D^{\gamma}_{\beta}\underbrace{ g_{\gamma \xi}g^{\xi\gamma} }_{ =1 }=-b_{\beta \xi}g^{\xi\gamma}
\label{e87292}
\end{equation}

\cref{e87292} 证明的逻辑是先任给指标 $\sigma$，对 \cref{cb4036} 左右乘上 $g^{\xi\sigma}$，自然成立
\[
(D^{\gamma}_{\beta}g_{\gamma \xi})g^{\xi\sigma}=-b_{\beta \xi}g^{\xi\sigma}
\]
再交换求和次序得到
\[
D^{\gamma}_{\beta}\underbrace{ (g_{\gamma \xi}g^{\xi\sigma}) }_{ =\delta_{\gamma}^{\sigma} }=-b_{\beta \xi}g^{\xi\sigma}
\]
令 $\sigma=\gamma$ 即可.

引入指标上升/下降的概念，命
\[
b^{\gamma}_{\beta}\coloneqq b_{\beta \xi}g^{\xi\gamma}
\]
从而
\[
b_{\beta\gamma}=b_{\beta}^{\xi}g_{\xi\gamma}
\]
接下来我们求 $\Gamma^{\gamma}_{\alpha\beta}$. 对于 \cref{1233fd} 第一式点乘 $r_{\xi}$ 得到
\[
r_{\alpha\beta}\cdot r_{\xi}=\Gamma^{\gamma}_{\alpha \beta}g_{\gamma \xi} \overset{ \text{指标下降} }{ = } \Gamma_{\xi\alpha\beta}
\]
注意到指标的对称性，以及 $g_{\alpha\beta}=r_{\alpha}\cdot r_{\beta}$，故
\[
\begin{cases}
\frac{ \partial g_{\alpha\beta} }{ \partial u^{\gamma} } & =r_{\alpha\gamma}\cdot r_{\gamma}+r_{\alpha}\cdot r_{\beta\gamma}=\Gamma_{\beta\alpha\gamma}+\Gamma_{\alpha\beta\gamma} \\
 \frac{ \partial g _{\alpha\gamma}}{ \partial u^{\beta} }  & =\Gamma_{\alpha\gamma\beta}+\Gamma_{\gamma\alpha\beta} \\
 \frac{ \partial g_{\gamma\beta} }{ \partial u^{\alpha} }  & =\Gamma_{\gamma\beta\alpha}+\Gamma_{\beta\gamma\alpha}
\end{cases}
\]
解得
\[
2\Gamma_{\gamma\alpha\beta}=\frac{ \partial g_{\alpha\gamma} }{ \partial u^{\beta} } +\frac{ \partial g_{\gamma\beta} }{ \partial u^{\alpha} } -\frac{ \partial g_{\alpha\beta} }{ \partial u^{\gamma} } 
\]
于是
\[
\Gamma^{\gamma}_{\alpha\beta}=\Gamma _{\xi\alpha\beta}g^{\gamma \xi}=\frac{1}{2}g^{\gamma \xi}\left( \frac{ \partial g_{\alpha \xi} }{ \partial u^{\beta} } +\frac{ \partial g_{\xi\beta} }{ \partial u^{\alpha} } -\frac{ \partial g_{\alpha\beta} }{ \partial u^{\xi} }  \right)
\]
综上，\cref{1233fd} 变为
\begin{equation}
\begin{cases}
\frac{ \partial r_{\alpha} }{ \partial u^{\beta} }  & =\Gamma^{\gamma}_{\alpha\beta}r_{\gamma}+b_{\alpha\beta}n \\
\frac{ \partial n }{ \partial u^{\beta} }  & =-b^{\gamma}_{\beta}r_{\gamma }
\end{cases}
\label{9c4151}
\end{equation}

这里 $\Gamma^{\gamma}_{\alpha\beta}$ 称为 \textbf{Christoffel 记号}. 通常把 \cref{9c4151} 的第一式称为 \textbf{Gauss 公式}，第二式称为 \textbf{Weingarten 公式}.

\begin{note}
$\Gamma^{\gamma}_{\alpha\beta}$ 的几何意义是 $r_{\alpha\beta}$ 用自然标架分解时在切向量 $r_{\gamma}$ 上的投影（是那个系数）；$\Gamma_{\gamma\alpha\beta}$ 的几何意义是 $r_{\alpha\beta}$ 在切向量 $r_{\gamma}$ 上的投影，即 $r_{\alpha \beta}\cdot r_{\gamma}$. 当 $r_1\not\perp r_2$ 时他们不同.
\end{note}
\begin{figure}[H]
\centering
\includegraphics[width=\textwidth]{differential-geometry-2025071000.png}
% \caption{}
\label{}
\end{figure}

\subsection{Riemann 符号}

若 $r\in C^{3}$，则
\[
r_{\alpha\beta\gamma}=r_{\alpha\gamma\beta},\qquad n_{\beta\gamma}=n_{\gamma\beta}
\]
由 \cref{9c4151} 可知
\begin{equation}
(\Gamma^{\sigma}_{\alpha\beta}r_{\delta}+b_{\alpha\beta}n)_{\gamma}=(\Gamma^{\delta}_{\alpha\gamma}r_{\delta}+b_{\alpha\gamma}n)_{\beta}
\label{djl1jd}
\end{equation}
\[
(b^{\delta}_{\beta}r_{\delta})_{\gamma}=(b_{\gamma}^{\delta}r_{\delta})_{\beta}
\]
展开得到并利用 \cref{9c4151} 化简得到两个式子
\[
(b_{\alpha\beta})_{\gamma}-(b_{\alpha\gamma})_{\beta}=b_{\gamma\delta}\Gamma^{\delta}_{\alpha\gamma}-b_{\gamma\delta}\Gamma^{\delta}_{\alpha\beta}
\]
\[
(\Gamma^{\delta}_{\alpha\beta})_{\gamma}-(\Gamma^{\delta}_{\alpha\gamma})_{\beta}+\Gamma^{\eta}_{\alpha\beta}\Gamma^{\delta}_{\eta\gamma}-\Gamma^{\eta}_{\alpha\gamma}\Gamma^{\delta}_{\eta\beta}=b_{\alpha\beta}b^{\delta}_{\gamma}-b_{\alpha\gamma}b^{\delta}_{\beta}\eqqcolon -R^{\delta}_{\alpha\beta\gamma}
\]
称为第一类基本量的 \textbf{Riemann 符号}.

\begin{note}
Riemann 符号的定义可能偏差一个负号.
\end{note}
同时，记指标下降为
\begin{equation}
R_{\delta\alpha\beta\gamma}\coloneqq g_{\delta \eta}R^{\eta}_{\alpha\beta\gamma}=b_{\delta\beta}b_{\alpha\gamma}-b_{\delta\gamma}b_{\alpha\beta}
\label{djl1jd2}
\end{equation}

\cref{djl1jd2} 称为曲面的 \textbf{Gauss 方程}， \cref{djl1jd} 称为 \textbf{Codazzi 方程}.

要指出 Riemann 符号具有对称性质：
\[
R_{\delta\alpha\beta\gamma}=R_{\beta\gamma\delta\alpha}=-R_{\alpha\delta\beta\gamma}=-R_{\delta\alpha\gamma\beta}
\]
本质上 \cref{djl1jd2} 只包含一个方程
\[
R_{1212}=b_{11}b_{22}-(b_{12})^{2}
\]
\cref{djl1jd} 中有意义的方程只有两个（若 $\gamma=\beta$ 则恒等）
\[
\begin{cases}
(b_{11})_{2}-(b_{12})_{1} & =-b_{2\gamma}\Gamma^{\gamma}_{11}+b_{1\gamma}\Gamma_{12}^{\gamma} \\
(b_{21})_{2}-(b_{22})_{1} & =-b_{2\gamma}\Gamma_{21} ^{\gamma}+b_{1\gamma}\Gamma_{22}^{\gamma}
\end{cases}
\]
\begin{figure}[H]
\centering
\includegraphics[width=\textwidth]{1-differential-geometry-2025071000.png}
% \caption{}
\label{}
\end{figure}
\begin{figure}[H]
\centering
\includegraphics[width=\textwidth]{2-differential-geometry-2025071000.png}
% \caption{}
\label{}
\end{figure}
\begin{figure}[H]
\centering
\includegraphics[width=\textwidth]{3-differential-geometry-2025071000.png}
% \caption{}
\label{}
\end{figure}

\subsection{唯一性、存在性}

\begin{figure}[H]
\centering
\includegraphics[width=\textwidth]{4-differential-geometry-2025071000.png}
% \caption{}
\label{}
\end{figure}

\subsection{Gauss 绝妙定理}

\begin{figure}[H]
\centering
\includegraphics[width=\textwidth]{5-differential-geometry-2025071000.png}
% \caption{}
\label{}
\end{figure}
\begin{figure}[H]
\centering
\includegraphics[width=\textwidth]{6-differential-geometry-2025071000.png}
% \caption{}
\label{}
\end{figure}

Gauss 绝妙定理启发我们对于只给定第一基本形式的曲面进行研究，这样的几何学称为内蕴几何学，它在高维的推广即 Riemann 几何学.


\end{document}